% Generated draft: requires mathematical and source review.
\documentclass[11pt]{article}
\usepackage[margin=25mm]{geometry}
\usepackage{fontspec}
\setmainfont{Latin Modern Roman}
\usepackage{amsmath,amssymb,mathtools,unicode-math}
\setmathfont{Latin Modern Math}
\usepackage{enumitem,needspace,xurl,hyperref}
\setlength{\parindent}{0pt}
\setlength{\parskip}{5pt}
\setlength{\emergencystretch}{3em}
\newcommand{\N}{\mathbb{N}}
\newcommand{\Z}{\mathbb{Z}}
\newcommand{\Q}{\mathbb{Q}}
\newcommand{\R}{\mathbb{R}}
\newcommand{\C}{\mathbb{C}}
\newcommand{\F}{\mathbb{F}}
\newcommand{\A}{\mathbb{A}}
\newcommand{\PP}{\mathbb{P}}
\newcommand{\E}{\mathbb{E}}
\newcommand{\eps}{\varepsilon}
\newcommand{\dd}{\,\mathrm{d}}
\begin{document}
\section*{UnsolvedMath ID: 701742}
\textbf{Source problem code:} catalogue entry
\subsection{Definitions and mathematical statement}
For a group \(G\), define its lower central series by \(\gamma_1(G)=G\) and
\[
\gamma_{n+1}(G)=\langle [x,y]:x\in\gamma_n(G),\ y\in G\rangle,
\qquad [x,y]=x^{-1}y^{-1}xy.
\]
The group \(G\) has nilpotency class exactly \(3\) if \(\gamma_4(G)=\{1\}\) and \(\gamma_3(G)\ne\{1\}\).

Under the conventional integral interpretation of dimension subgroups, let \(\Z[G]\) be the integral group ring, let
\[
\varepsilon:\Z[G]\longrightarrow\Z,
\qquad \varepsilon\!\left(\sum_{g\in G}a_g g\right)=\sum_{g\in G}a_g,
\]
where each sum has finite support, and let \(I(G)=\ker\varepsilon\) be the augmentation ideal. Define
\[
D_5^{\Z}(G)=\{g\in G:g-1\in I(G)^5\},
\]
where \(I(G)^5\) is the fifth ideal power.

With this interpretation, the question proposed by E. Rips is whether
\[
\exists\text{ a group }G\quad
\gamma_4(G)=\{1\},\qquad
\gamma_3(G)\ne\{1\},\qquad
D_5^{\Z}(G)\ne\{1\}.
\]
No finiteness or finite generation hypothesis is imposed in the supplied statement. The coefficient ring for the dimension subgroup is not explicitly specified there; the integral interpretation above requires confirmation against the original source.
\subsection{Short English statement}
Does there exist a group of nilpotency class exactly \(3\) whose fifth dimension subgroup is nontrivial? Under the conventional integral interpretation, this means that some \(g\ne1\) satisfies \(g-1\in I(G)^5\) in \(\Z[G]\). The supplied record does not explicitly specify the coefficient ring.
\subsection{Sources}
\begin{itemize}
\item Local UnsolvedMath registry, permanent ID \(701742\), exact source problem code \texttt{catalogue entry}, titled \emph{Kourovka Notebook Problem 20.84}. The supplied registry labels the problem open; this label has not been independently verified.
\item \emph{The Kourovka Notebook (Unsolved Problems in Group Theory)}, original problem \(20.84\), proposed by E. Rips in \(2022\). The supplied background identifies the 20th issue (2022) and the 21st issue edition (2026), \texttt{arXiv:1401.0300v48}. The supplied collection description names E. I. Khukhro and V. D. Mazurov as editors.
\item The supplied source text and bibliographic attribution were used without independent consultation of the cited editions.
\end{itemize}
\end{document}
